\section{Background}
\label{sec:background}

\subsection{Congruence closure}

The \emph{congruence closure} procedure decides whether a finite conjunction of
equalities and disequalities over uninterpreted function symbols is satisfiable.
Two terms are \emph{congruent} when they have the same operator and
pairwise-congruent children: $f(s_1, \ldots, s_k) \equiv f(t_1, \ldots, t_k)$
whenever $s_i \equiv t_i$ for every $i$. Together with reflexivity, symmetry,
and transitivity of equality, this gives an equivalence relation on terms;
closing under it propagates a small set of asserted equalities to all of their
consequences.

If an SMT solver is given a disequality $s \neq t$ and if $s$ and $t$ are in the
same equivalence class, then it knows the formula is unsatisfiable. Similarly,
an equality saturation engine may ask if two programs $P$ and $Q$ are equivalent
under a set of rewrite rules by checking if $P$ and $Q$ are in the same
equivalence class.

Hence, the decision problem of congruence closure is to determine
whether any two terms are in the same equivalence class
after closing under the asserted equalities and congruence.

The classical algorithm of Nelson and
Oppen~\cite{nelson-oppen-cc} maintains a union-find over equivalence
classes of subterms, plus a hashcons that maps every canonicalized
e-node to its current class. A worklist holds classes that have just
been merged. Processing a class involves walking its
\emph{parent list}---all e-nodes that mention this class as a direct
child---re-canonicalizing each, and merging any pair that the hashcons
identifies as a duplicate. New merges produce new worklist entries; the
loop runs to a fixed point. 

\subsection{Motivating Example}
\label{sec:bg:motivating}

To make the structure concrete, consider the smallest non-trivial
instance of what we call the \emph{quadratic} family of benchmarks
(Section~\ref{sec:eval}). For parameter $n = 2$:
%
\begin{verbatim}
(declare-fun f (U U) U)
(declare-fun g (U U) U)
(declare-const a0 U)  (declare-const b0 U)
(declare-const a1 U)  (declare-const b1 U)
(assert (= a0 b0))
(assert (= a1 b1))
(assert (not (=
  (g (g (f a0 a0) (f a0 a1))
     (g (f a1 a0) (f a1 a1)))
  (g (g (f b0 b0) (f b0 b1))
     (g (f b1 b0) (f b1 b1))))))
(check-sat)
\end{verbatim}
%
There are two asserted equalities, $a_0 = b_0$ and $a_1 = b_1$, and a
single disequality between two terms built from $f$ and $g$. The
disequality compares two balanced binary trees of $g$, each whose
leaves are the four $f$-terms over the index set $\{0, 1\}^2$. The
expected outcome is \texttt{unsat}.

Closure proceeds as follows. After applying the two asserted unions,
the union-find places $a_0 \equiv b_0$ and $a_1 \equiv b_1$. The
\emph{f-layer} now sees that
$f(a_0, a_0) \equiv f(b_0, b_0)$ by congruence (both children
match), and similarly for the other three index pairs. That is
$n^2 = 4$ congruence merges in a single step. Once those fire, the
$g$-tree's bottom layer becomes pairwise congruent: $g(f(a_0,a_0),
f(a_0,a_1)) \equiv g(f(b_0,b_0), f(b_0,b_1))$, and the analogous
right-half merge. One step higher and the two roots collapse, at
which point the disequality is violated.

Two features make this a clean stress test for parallelism:

\begin{itemize}
\item \textbf{All $n^k$ f-layer congruences are mutually independent.}
      Their canonical signatures depend only on the post-merge roots
      of the leaves, which are fixed once the asserted unions are
      applied. A bulk-synchronous closure pass can detect all of
      them in a single round, with no inter-merge serialization.

\item \textbf{The remaining work is a logarithmic dependency chain.}
      Climbing the $g$-tree requires one round per layer, and each
      layer halves in size. So the total round count is
      $1 + \lceil \log_2(n^k) \rceil$: one for the f-layer, and
      $\lceil \log_2(n^k) \rceil$ for the $g$-tree.
\end{itemize}

\noindent
We make these intuitions precise in
Section~\ref{sec:alg}. Increasing $k$ (the arity of $f$) changes the
benchmark from quadratic ($n^2$ congruences) to cubic ($n^3$),
quartic ($n^4$), or quintic ($n^5$), allowing us to dial up the work
density per round without changing the round count by much.

\subsection{ParlayLib and BSP-style closure}

Our implementation uses ParlayLib~\cite{parlaylib}, a C++ work-stealing
runtime with composable parallel primitives over a contiguous
\texttt{parlay::sequence} type. The primitives we rely on most
directly are \texttt{parallel\_for}, \texttt{filter},
\texttt{flatten}, and a parallel \emph{semisort}~\cite{semisort}. The
semisort takes a sequence of (hash, payload) pairs and reorganizes
them so that all entries with the same hash are contiguous; we use it
to group e-nodes by canonical signature within a closure round.
ParlayLib's scheduler controls the number of worker threads via the
\texttt{PARLAY\_NUM\_THREADS} environment variable, which we use to
sweep core counts in our evaluation.

A \emph{bulk-synchronous parallel} (BSP) pass alternates between
parallel work and global barriers. Our parallel closure is structured
this way: each round runs a full data-parallel pipeline (gather
frontier $\to$ canonicalize $\to$ semisort by signature $\to$ merge
groups), then the next round starts only once the previous one's
unions are visible to all workers. Within a round, every step of the
pipeline is a single parlay primitive over a single sequence, so the
runtime can stripe the work across threads without us writing any
explicit synchronization.